\documentclass[12pt,a4paper]{article}
\usepackage{amsmath,amssymb,amsthm}
\usepackage{graphicx}
\usepackage{geometry}
\geometry{a4paper, left=2.5cm, right=3cm, top=3cm, bottom=3cm}
\usepackage{fancyhdr}
\usepackage{setspace}
\usepackage{booktabs}
\usepackage{float}
\usepackage{hyperref}
\usepackage{xepersian}

\settextfont{X Nazanin}
\setdigitfont{X Nazanin}
\setlatintextfont{Times New Roman}

\makeatletter
\def\UseMathForPositioningText{}
\def\@kernel@tabular@init{}
\makeatother

\begin{document}

\begin{center}
{\Large\bfseries دانشگاه تهران}\\[0.2cm]
{\large\bfseries دانشکده مهندسی برق و کامپیوتر}\\[0.2cm]
{\large\bfseries گزارش فنی و تحلیل معماری مکانیزم توجه}\\[0.5cm]
\rule{\linewidth}{1pt}\\[0.3cm]
{\LARGE\bfseries ماژول توجه: SOAFusedAttention}\\[0.2cm]
{\large \textbf{شناسه رسمی در فریم‌ورک:} \lr{SOA-Fused Attention: Low-Rank Latent Space Multi-Head Attention with Zero-Shot SVD Initialization}}\\[0.2cm]
\rule{\linewidth}{1pt}\\[0.5cm]
\end{center}

\begin{figure}[H]
\centering
\framebox[0.9\textwidth]{
\begin{minipage}{0.85\textwidth}
\vspace{0.3cm}
\centering
\textbf{[محل الحاق دیاگرام تخصصی این مکانیزم]}\\[0.2cm]
\small \textbf{پرامپت مهندسی‌شده هوش مصنوعی جهت ساخت تصویر:}\\[0.15cm]
\footnotesize \lr{A detailed, high-tech academic architecture diagram of SOAFusedAttention (SOA-Fused Attention: Low-Rank Latent Space Multi-Head Attention with Zero-Shot SVD Initialization). Visualizing Query, Key, Value tensor transformations, matrix operations, and memory optimizations inside GPU hardware cache. Glowing cybernetic blueprint style, dark background, cyan and deep amber highlights.}
\vspace{0.3cm}
\end{minipage}
}
\caption{دیاگرام مفهومی و شماتیک جریان داده در مکانیزم توجه SOAFusedAttention}
\end{figure}

\section*{۱. کارنامه تجربی و شاخص‌های بنچ‌مارک}
\begin{itemize}
\item \textbf{وضعیت ارزیابی در خط لوله:} \lr{PASS}
\item \textbf{دقت ثبت‌شده روی بنچ‌مارک:} 81.31 درصد
\item \textbf{امتیاز \lr{F1}:} 0.8191
\item \textbf{اوج مصرف حافظه فعال‌سازی (\lr{Peak VRAM}):} 282.69 مگابایت
\end{itemize}
\textbf{تحلیل و توجیه علمی نتایج:} این ماژول با ثبت دقت 81.31 درصد مصالحه مطلوبی میان سرعت اجرای استنتاج و بازنمایی ویژگی‌های توکن‌ها به نمایش می‌گذارد. تقریب‌های اعمال‌شده باعث مهار گلوگاه پهنای باند حافظه گردیده‌اند.

\section*{۲. مبانی تئوریک و انگیزه طراحی معماری}
Dense Multi-Head Self-Attention layers allocate distinct projection matrices for Query, Key, and Value ( W Q, W K, W V in mathbb R D times D ). In standard architectures, these matrices contain significant spectral redundancy, with >80 of their Frobenius energy concentrated in low-rank singular subspaces.

\section*{۳. روابط و فرمول‌بندی دقیق ریاضی}
2.1 Low-Rank Factorization of Attention Projections

\section*{۴. تحلیل پیچیدگی زمانی و رفتار حافظه سخت‌افزاری}
---

\end{document}
